% !TeX program = lualatex
% !TeX encoding = utf8
% !TeX spellcheck = uk_UA
% !BIB program = biber

\documentclass[biblatex]{ConspectBook}

%===============================================================================
%
%							Вихідні дані докумета
%
%===============================================================================

\title{Елементи фізики плазми}
\def\subtitle{}
\def\authors{С.~М.~Пономаренко}

%=================================== Бібліографія  =============================
%

\addbibresource{./bibliography/bibliography.bib}

%
%===============================================================================

%===============================================================================
%
% !TeX root =./ClassicalElectrodynamics.tex

% chktex-file 1

%========================================================================================================
%
%									      Протоколи
%
%========================================================================================================


% ============ НН ФТІ ==========
\newcommand\IPTNo{XX}            % Номер протокола затвердження методичною комісією НН ФТІ
\newcommand\IPTDate{XX.XX.2026}  % Дата протокола затвердження методичною комісією НН ФТІ

% ============ КПІ =============
\newcommand\KPINo{X}             % Номер протокола надання грифу КПІ
\newcommand\KPIDate{XX.XX.2026}  % Дата протокола надання грифу КПІ

%========================================================================================================
%
%									      Палітурка
%
%========================================================================================================

\newcommand{\CoverPage}{
	\begin{alwayssingle}
		\begin{center}
			\begin{flushright}\bfseries\sffamily
				\MakeUppercase{Міністерство освіти і науки України}\\
				КПІ ім. Ігоря Сікорського\\
			\end{flushright}
			\begin{tcolorbox}[titlepagestyle,
					toprule=0.10cm,
					bottomrule=0.10cm,
					overlay={%
						\node (picture) at ([xshift=4cm]frame.west) {\includegraphics[width=3.5cm]{logo/logo_IPT}};
					}
			]%
			\begin{flushright}
				\large\bfseries\color{white}Навчально-науковий\\ фізико-технічний інститут
			\end{flushright}
			\end{tcolorbox}
			\vspace*{-2em}
		 	\begin{Large}\color{themecolordark!90!black}
			\begin{alignat*}{3}
				\vect{\nabla}&\times&&\Efield &&= -\dfrac{1}{c}\dfrac{\partial\Bfield}{\partial t} \\
				\vect{\nabla}&\,\cdot&&\Bfield &&= 0 \\
				\vect{\nabla}&\,\cdot&&\Dfield &&= 4\pi\rho \\
				\vect{\nabla}&\times&&\Hfield &&= \dfrac{4\pi}{c} \vect{j}+\dfrac{1}{c}\dfrac{\partial\Dfield}{\partial t}
			\end{alignat*}
			\end{Large}
			\vspace*{2em}
			\begin{tcolorbox}[
				titlepagestyle,
				toprule=0.15cm,
				bottomrule=0.15cm,
				top=1.3cm,
				bottom=0.7cm,
				overlay={%
				\node[%
							fill=white,
							rounded corners = 15pt,
							draw=themecolorlight,
							line width=0.15cm,
							inner sep=0pt,
							text width=17cm,
							minimum height=2cm,
							align=center,
							%anchor=east,
							font=\sffamily\bfseries\large
						] (title) at (frame.north) {С.~М. Пономаренко
					};
				}
			]
			\centering
			\Huge\sffamily\bfseries\textcolor{white}{\realtitle}\\
			\huge\sffamily\bfseries\textcolor{white}{\subtitle}
			\end{tcolorbox}
			\vfill
		 	\input{./MaxwellEqns/tikz/cover.tikz}
			\vfill
			\begin{tcolorbox}[titlepagestyle,
					toprule=0.10cm,
					bottomrule=0.10cm]
				\begin{center}\color{white}\bfseries\normalsize
					\MakeUppercase{Київ~\the\year} \\
%					КПІ ім. Ігоря Сікорського \\
%					\the\year
				\end{center}
			\end{tcolorbox}
		\end{center}
		\clearpage
	\end{alwayssingle}
\setcounter{page}{1}
}


%========================================================================================================
%
%									      Титульна сторінка
%
%========================================================================================================

\renewcommand\maketitle{
	\begin{alwayssingle}
		\begin{center}
				\MakeUppercase{Міністерство освіти і науки України}

				\bigskip
				\MakeUppercase{Національний технічний університет України}\\
				<<КИЇВСЬКИЙ ПОЛІТЕХНІЧНИЙ ІНСТИТУТ \\ імені ІГОРЯ СІКОРСЬКОГО>>
				\vspace*{100pt}

				{\large С.~М.~Пономаренко
						}
				\vspace*{50pt}

				{\Huge\sffamily\bfseries\realtitle}\\[1em]
				{\huge\sffamily\bfseries\subtitle}

			\vspace*{50pt}
			\begin{center}\itshape
			Рекомендовано Методичною радою КПІ ім. Ігоря Сікорського як навчальний посібник для
			здобувачів ступеня бакалавра за спеціальностями E6 Прикладна фізика та наноматеріали.
			\end{center}
            \vfill
%                 \textcolor{malina}{Видання друге, виправлене та доповнене}
%                \textcolor{malina}{Версія від~\today, час компіляції \currenttime.}
			\vfill
			\begin{center}
				\MakeUppercase{Київ} \\
				КПІ ім. Ігоря Сікорського \\
				\the\year
			\end{center}
		\end{center}
		\clearpage
	\end{alwayssingle}
}


%========================================================================================================
%
%									      Друга сторінка
%
%========================================================================================================
\newcommand\makeinfopage{
	\begin{alwayssingle}
		\noindent%
		\begin{minipage}[t]{0.5\textwidth}
				\begin{flushleft}
					УДК  537(075.8)\\
					ББК  22.31я73\\
					П 563
				\end{flushleft}
		\end{minipage}


		\bigskip\noindent%
        \begin{minipage}[t]{0.25\linewidth}
            	\begin{flushleft}
                    Рецензенти:
                \end{flushleft}
        \end{minipage}\hfill
        \begin{minipage}[t]{0.75\linewidth}
%                \href{}{}, д.ф.-м.н., професор кафедри загальної фізики фізико-метематичного факультету КПІ ім. Ігоря Сікорського.
%                \\[1ex]
%                \href{https://phys.knu.ua/1450-vilchinskiy-stanislav-yosipovich}{С.~Й.~Вільчинський},  д.ф.-м.н., професор, завідувач кафедри  квантової теорії поля та космомікрофізики.
        \end{minipage}

		\bigskip\noindent%
        \begin{minipage}[t]{0.25\linewidth}
            	\begin{flushleft}
                    Відповідальний редактор:
                \end{flushleft}
        \end{minipage}\hfill
        \begin{minipage}[t]{0.75\linewidth}
                \href{http://is.ipt.kpi.ua/is/smirnovsa/}{С.~А.~Смирнов}, к.ф.-м.н., доцент, голова
                методичної ради НН ФТІ КПІ ім. Ігоря Сікорського
        \end{minipage}

		\begin{center}\itshape\small
				Гриф надано Методичною радою КПІ ім. Ігоря Сікорського (протокол №~\KPINo~від \KPIDate~р.) за поданням Вченої ради Навчально-наукового Фізико-технічного інституту (протокол №~\IPTNo~від \IPTDate~р.)
		\end{center}
		\begin{center}
			\ifelectronic Електронне мережне навчальне видання \fi
			%\par {Версія від~\href{http://www.istpravda.com.ua/dates}{\today}} \par\else \par  \fi
		\end{center}
		\begin{center}
			\href{https://apd.ipt.kpi.ua/kafedra/stuff/ponomarenko}{\itshape Пономаренко Сергій
			Миколайович}, к.ф.-м.н., доцент
		\end{center}
%			\vspace*{1em}%
		\begin{center}\bfseries
			\LARGE\sffamily\realtitle \\
			\Large\sffamily\subtitle
		\end{center}
        \noindent%
        \begin{minipage}[t]{\textwidth}\small
                \realtitle: \subtitle\ [Електронний ресурс] : навч. посіб. для студ. спеціальностей
                E6 <<Прикладна фізика та наноматеріали>> / С.~М. Пономаренко; КПІ ім. Ігоря Сікорського.~--- Електронні
                текстові дані
            (1 файл: 789~кБ). – Київ: КПІ ім. Ігоря Сікорського, \the\year. ---
             \the\numexpr\getpagerefnumber{LastPages}-1\relax~с.
        \end{minipage}

%		\noindent%
%		\begin{flushleft}
%			\begin{tabular}{lp{0.9\textwidth}}
%				     & В.~І. Жданов, С.~М. Пономаренко, В.~Б. Долгошей                                                                                                                                                                        \\
%				Ж 42 & \hspace*{3ex} \realtitle : \subtitle{} [Електронний ресурс] : навчальний посібник / %
%				В.~І. Жданов, %
%				С.~М. Пономаренко%,
%				В.~Б. Долгошей %
%				--- К.:~КПІ ім. Ігоря Сікорського, \the\year. --~\the\numexpr\getpagerefnumber{LastPage}-1\relax~с. -- Бібліогр.: с.~\pageref{BibPage}\relax. \ifelectronic\relax\else-- 80~прим.\fi
%			\end{tabular}
%		\end{flushleft}
		\vfill


		Для студентів фізико-технічного інституту КПІ ім. Ігоря Сікорського, які навчаються за спеціальностями E6~Прикладна фізика та наноматеріали.

		\vfill

%		\begin{flushleft}\small
%			Ілюстративний матеріал підручника підготовлений за допомогою пакету \href{http://pgf.sourceforge.net}{TikZ/Pgf}. Верстка тексту проведена в видавничій системі \LaTeXe{} (компілятор Lua\LaTeX) на базі системи комп'ютерної верстки \TeX{} (Збірка  \href{https://www.tug.org/texlive/}{\TeX Live~\the\year}) з використанням оболонки \href{https://www.texstudio.org}{\TeX Studio}.
%		\end{flushleft}
	\hfill
        \begin{flushright}
    	\begin{minipage}[t]{0.6\linewidth}\small
        \textcopyright{} С.~М. Пономаренко, \the\year\,р. \\
        \textcopyright{}  КПІ ім. Ігоря Сікорського (НН ФТІ), \the\year~р.
        \end{minipage}
        \end{flushright}
		\newpage%
	\end{alwayssingle}
}

%%%%%%%%%%%%%%%%%%%%%%%%%%%%%%%%%%%%%%%%%%%%%%%%%%%%%%%%%%%%%%%%%%%%%%%%%%%%
%%                                                                        %%
%%                              Last Page                                 %%
%%                                                                        %%
%%%%%%%%%%%%%%%%%%%%%%%%%%%%%%%%%%%%%%%%%%%%%%%%%%%%%%%%%%%%%%%%%%%%%%%%%%%%
\newcommand{\makelastpage}{%
\clearpage%
\thispagestyle{empty}%
\vspace*{0.4\textheight}
\begin{center}
	    \textbf{Пономаренко} Сергій Миколайович
\end{center}

\begin{center}\bfseries
    \Large\sffamily\realtitle \\
    \large\sffamily\subtitle
\end{center}

\begin{center}\itshape
    Комп'ютерне верстання в системі \LaTeXe{} С.\ М. Пономаренко
\end{center}

\vspace*{1em}
\begin{center}\small
Національний технічний університет України \\
<<Київський політехнічний інститут імені Ігоря Сікорського>> \\
Свідоцтво про державну реєстрацію: серія ДК № 5354 від 25.05.2017 р.\\
просп. Берестейський, 37, м. Київ, 03056
\end{center}

%\vfill
%
%\begin{center}
%    Видавництво <<Політехніка>> \kpishort \\
%    вул. Політехнічна, 14, корп. 15 \\
%    м. Київ, 03056 \\
%    Тел. (044) 204-81-78
%\end{center}
}


%			Титульна сторінка
%
%===============================================================================
\begin{document}
\pagestyle{empty}
%\ifelectronic\CoverPage\setcounter{page}{0}\fi
% \CoverPage
\maketitle
\makeinfopage
%===============================================================================
%
\clearpage\pagestyle{plain}
\phantomsection
\addtocontents{toc}{\protect\vfill}
\addcontentsline{toc}{chapter}{Зміст}
\tableofcontents%				   Зміст
%
%===============================================================================

%===============================================================================
%
%\include{intro}%                  Вступ та передмова
%
%===============================================================================

%===============================================================================
%
%									      Вставка файлів розділів
%
%===============================================================================
%\setcounter{chapter}{-1}
%\input{Additions/MaxwellEquation}

%===============================================================================

\pagestyle{main}
% !TeX program = lualatex
% !TeX encoding = utf8
% !TeX spellcheck = uk_UA
% !TEX indexprogram = upmendex
% !TeX root =./main.tex

%=========================================================
\chapter*{Передмова}\label{\currfilebase}\hypertarget{\currfilebase}{}
\addcontentsline{toc}{chapter}{Передмова}
%=========================================================


Переважна більшість баріонної речовини у~Всесвіті~--- зорі, міжзоряний газ, сонячний
вітер~--- перебуває в~стані плазми. На Землі плазма~--- це блискавка, полярні сяйва,
іоносфера, газовий розряд у~лампах, зварювальна дуга, плазмове травлення мікросхем
і~плазма, яку намагаються утримувати в~установках керованого термоядерного синтезу.

Спільна риса всіх цих явищ~--- колективна поведінка заряджених частинок. У~всьому розділі
важливо відрізняти \emphx{масштабну оцінку} від точної теорії. Якщо результат отримано
в~простій моделі, ми спочатку з'ясуємо фізичний зміст цієї моделі, а~потім зазначимо, де
вона перестає працювати.




Фізика плазми поєднує кілька рівнів опису. Спочатку ми
визначимо, коли система заряджених частинок справді поводиться як плазма, і~введемо її
характерні просторові та часові масштаби. Потім розглянемо рух \emph{окремої} частинки
в~електромагнітному полі, зіткнення та перенос. Після цього перейдемо до
\emph{колективних} коливань і~хвиль, а~далі --- до гідродинамічного та
магнітогідродинамічного опису. Нарешті застосуємо ці моделі до магнітного утримання
і~керованого термоядерного синтезу.


\emphx{Головний принцип посібника:} одна й~та сама плазма може потребувати різних моделей
залежно від масштабу явища. Тому щоразу будемо ставити три питання: \emph{який просторовий
масштаб? який часовий масштаб? які фізичні процеси встигають відбутися?}


%% --------------------------------------------------------
\section*{Система одиниць}
%% --------------------------------------------------------

У~цьому посібнику, як і~в~більшості наведених нижче формул, використано систему СГС-Гаусса:
сила Лоренца має вигляд $\vect F=q(\vect E+\vect v\times\vect B/c)$, а~закон Ампера ---
$\rot\vect B=(4\pi/c)\vect j+(1/c)\partial_t\vect E$. Наприкінці розділу наведено
практичні формули в~СГС. Перехід до СІ простий, але множники $c$ та $4\pi$ в~окремих
рівняннях змінюються.

\multiinclude{
    Basics,
    Waves,
    Motion,
    Fluid,
    LinearWaves,
    Kinetics,
    Collisions,
    Instability,
    Confinement,
    Fusion,
    Scales,
}

\chapter*{Підсумки}
\phantomsection\addcontentsline{toc}{chapter}{Підсумки}
%% --------------------------------------------------------

\begin{itemize}
	\item Іонізація сама по собі ще не робить газ плазмою в~колективному сенсі. Потрібно,
	  щоб на дебаївській довжині $\lambda_D$ містилося багато частинок: $N_D\gg1$. Тоді
	  $a\ll\lambda_D$, а~на масштабах $L\gg\lambda_D$ плазма квазінейтральна.
	\item Дебаївське екранування дає просторовий масштаб колективної відповіді, а~плазмова
	  частота $\omega_p$~--- часовий. Для електронної компоненти
	  $\lambda_{De}=v_{Te}/\omega_{pe}$.
	\item Здатність електронів рухатися у~відповідь на змінне поле робить діелектричну
	  проникність плазми частотно-залежною: без зіткнень
	  $\varepsilon=1-\omega_p^2/\omega^2$. Звідси випливає частота відсікання та відбиття
	  низькочастотних хвиль.
	\item Зіткнення вводять дисипацію. Кулонівська частота $\nu_{ei}$ зазвичай значно
	  менша за $\omega_p$ у~слабкозв'язаній плазмі, але саме вона визначає опір,
	  теплоперенос і~дифузію.
	\item У~магнітному полі рух окремої частинки розпадається на швидкий ларморівський рух
	  і~повільні дрейфи. Найважливіші масштаби~--- $r_L$ та $\omega_c^{-1}$. Адіабатичний
	  інваріант $\mu=mv_\perp^2/(2B)$ пояснює дію магнітного дзеркала.
	\item На великих масштабах індивідуальний рух можна усереднити, і~плазма описується як
	  провідне середовище. У~МГД баланс газового та магнітного тисків має вигляд $\nabla
	  p=\vect j\times\vect B/c$, а~відношення цих тисків задає $\beta=8\pi p/B^2$.
	\item Магнітне поле може бути вморожене в~плазму, якщо дифузія поля повільна порівняно
	  з~його переносом: $R_m\gg1$. Але самого факту вмороженості недостатньо для
	  утримання~--- геометрія поля повинна ще забезпечувати замкнені та стійкі магнітні
	  поверхні.
	\item Магнітне утримання та керований синтез обмежуються нестійкостями й~переносом.
	  Тому межі простих моделей важливі не менше, ніж самі формули: вони показують, коли
	  потрібно переходити від частинкової картини до МГД, а~від МГД~--- до кінетичної
	  теорії.
	  	\item Силу, з~якою поле діє на плазму, можна записати як тиск $B^2/8\pi$ у~всіх
	  	  напрямках плюс натяг $B^2/4\pi$ уздовж ліній; через натяг пояснюються
	  	  альвенівські хвилі, обручева сила пінча та розпирання тороїдального кільця
	  	  зі~струмом.
	\item Частинкова й~рідинна картини сумісні: поперек поля гірація замінює зіткнення,
	  вмороженість є наслідком дрейфу $\vect E\times\vect B$, діамагнітний струм однаковий
	  у~частинковому, магнітному та~рідинному описах. Уздовж поля й~для тонких ефектів
	  потрібна кінетика.
	\item Енергетичний баланс плазми пов'язаний з~теоремою Пойнтінга; величина $\vect
	  j\cdot\vect E$ визначає локальну передачу енергії між полем і~частинками. Параметр
	  $\beta=8\pi p/B^2$ показує співвідношення теплового та~магнітного тисків.
\item Нестійкість означає зростання малого збурення рівноваги; її визначають знаком
  $\operatorname{Im}\omega$ після лінеаризації рівнянь. Джерелом вільної енергії можуть
  бути градієнти, струми або анізотропія функції розподілу.
\item Енергію синтезу переносять: $20\%$ несуть $\alpha$-частинки (лишаються в~плазмі),
  $80\%$~--- нейтрони (у~бланкеті дають тепло й~тритій). Плазмовий $Q$ не дорівнює
  інженерному: електростанції потрібні значно більші $Q$, ніж просте $Q>1$.
\end{itemize}

%%==============================================================================
%
%									       Вставка бібліографії
%
%===============================================================================
\clearpage\bibliographytrue\phantomsection\label{BibPage}
\pagestyle{bibliography}

\nocite{*}

\printbibliography[heading=bibintoc]

\makelastpage
\end{document}